% Generated draft: requires mathematical and source review.
\documentclass[11pt]{article}
\usepackage[margin=25mm]{geometry}
\usepackage{fontspec}
\setmainfont{Latin Modern Roman}
\usepackage{amsmath,amssymb,mathtools,unicode-math}
\setmathfont{Latin Modern Math}
\usepackage{enumitem,needspace,xurl,hyperref}
\setlength{\parindent}{0pt}
\setlength{\parskip}{5pt}
\setlength{\emergencystretch}{3em}
\newcommand{\N}{\mathbb{N}}
\newcommand{\Z}{\mathbb{Z}}
\newcommand{\Q}{\mathbb{Q}}
\newcommand{\R}{\mathbb{R}}
\newcommand{\C}{\mathbb{C}}
\newcommand{\F}{\mathbb{F}}
\newcommand{\A}{\mathbb{A}}
\newcommand{\PP}{\mathbb{P}}
\newcommand{\E}{\mathbb{E}}
\newcommand{\eps}{\varepsilon}
\newcommand{\dd}{\,\mathrm{d}}
\begin{document}
\section*{UnsolvedMath ID: 701100}
\textbf{Source problem code:} catalogue entry
\subsection{Definitions and mathematical statement}
Let \(\mathcal T\) be a rooted tree with root \(o\), and let \(V_n\) denote its vertices at distance \(n\) from \(o\), for \(n\in\Z_{\geq 0}\). For a vertex \(v\), let \(\mathcal T_v\) be the subtree consisting of \(v\) and its descendants. Suppose that a group \(G\) acts by root-preserving automorphisms on \(\mathcal T\). The action is faithful if the only element fixing every vertex is the identity, and is level-transitive if \(G\) acts transitively on each nonempty \(V_n\).

Define the rigid vertex stabilizer and the rigid level stabilizer by
\[
\operatorname{rist}_G(v)=\{g\in G: g(w)=w\text{ for every vertex }w\notin\mathcal T_v\},
\qquad
\operatorname{rist}_G(n)=\left\langle\operatorname{rist}_G(v):v\in V_n\right\rangle.
\]
In the terminology supplied by the source, \(G\) is a branch group for this action if the action is faithful and level-transitive and
\[
[G:\operatorname{rist}_G(n)]<\infty
\qquad\text{for every }n\in\Z_{\geq 0}.
\]

For an abstract group, interpret Kazhdan's property \((T)\) using the discrete topology: every unitary representation \(\pi:G\to\mathcal U(H)\) on a complex Hilbert space that has almost invariant unit vectors has a nonzero invariant vector. Here almost invariant unit vectors means that for every finite subset \(F\subseteq G\) and every \(\varepsilon>0\), there is \(\xi\in H\) with \(\|\xi\|=1\) and
\[
\|\pi(g)\xi-\xi\|<\varepsilon\qquad(g\in F).
\]
An invariant vector \(\eta\) satisfies \(\pi(g)\eta=\eta\) for every \(g\in G\).

\textbf{Question (supplied formulation).} Does there exist a group \(G\), together with a faithful, level-transitive action on a rooted tree \(\mathcal T\), such that every \(\operatorname{rist}_G(n)\) has finite index in \(G\) and \(G\) has Kazhdan's property \((T)\)?

\textbf{Unresolved qualifications.} The supplied record does not explicitly require the tree to be infinite, locally finite, spherically homogeneous, or without terminal vertices; it also gives no explicit infinitude or finite-generation hypothesis on \(G\). Such additional conventions can affect the intended existence question and must be checked before this draft is treated as a precise statement of the intended branch-group problem. No such hypotheses have been added here. The discrete interpretation of property \((T)\) is explicit in this draft; the record specifies no topology on \(G\).
\subsection{Short English statement}
The supplied question asks whether a group with Kazhdan's property \((T)\) can act faithfully and transitively on each level of a rooted tree while the subgroups generated by automorphisms supported inside individual subtrees have finite index at every level. The intended restrictions on the tree and group require clarification. The registry labels the problem open; that label is supplied source data, not an independently verified mathematical assessment.
\subsection{Sources}
\begin{itemize}
\item Supplied local UnsolvedMath registry record: permanent ID \(701100\), exact source code \texttt{catalogue entry}, title \emph{Kourovka Notebook Problem 15.12}, upstream revision \texttt{d371ee917ac9647fbb06cd6fc529deec6f3dda82}. The record was created and updated on 2026-10-16 and supplies the status \texttt{open}; these metadata do not establish an independently verified literature assessment.
\item \emph{The Kourovka Notebook (Unsolved Problems in Group Theory)}, problem 15.12, 15th issue (2002); the supplied background also identifies the 21st issue edition (2026), arXiv:1401.0300v48. The supplied collection description names E. I. Khukhro and V. D. Mazurov as editors. These bibliographic details are reproduced from the record and have not been independently checked.
\item Attribution supplied by the record: L. Bartholdi, R. I. Grigorchuk, and a third proposer whose name is corrupted in the supplied fields. The record gives the proposed year as 2002. The third name cannot be reliably reconstructed from the supplied text.
\item The statement cites R. I. Grigorchuk, in \emph{New horizons in pro-\(p\) groups}, Birkh\"auser, Boston, 2000, pp. 121--179, for motivation, examples, and known results. No contribution title is supplied.
\item The source includes the cross-reference ``See A14.34.'' Its target and contents are not included in the supplied record.
\end{itemize}
\end{document}
